\subsection*{Self-state representation}
\label{sec:model:state}

We model each bee agent as maintaining a four-dimensional self-state vector:
\begin{equation}
  \mathbf{s} = (\precision,\; c,\; \phi,\; e) \in [0,1]^4
\end{equation}
where $\precision$ is the \emph{precision} (inverse variance of internal predictions), $c$ is the
\emph{caste identity} self-representation ($0 =$ nurse, $1 =$ forager), $\phi$ is the
\emph{circadian phase accuracy}, and $e$ is the \emph{energy state}. The derived uncertainty is
$\uncertainty = 1 - \precision$.

Precision is the central parameter of the model. In the FEP framework~\cite{Friston2010},
precision weights prediction errors: high-precision agents trust their predictions and update beliefs
slowly; low-precision agents are driven by every sensory surprise. We implement precision as a scalar
for tractability, acknowledging that the biological CX likely employs multiple precision channels.

\subsection*{Empirical parameter grounding}
\label{sec:model:empirical}

We do not hand-tune precision distributions. Instead, we derive them from published bee psychophysics
data. Reversal learning studies~\cite{AvarWeber2012} report $d'$ values spanning roughly $[0.3, 2.1]$
across individual honeybees ($d'_{\max} \approx 2.5$ in optimal conditions). We normalize:
\begin{equation}
  \precision_i = 0.2 + \frac{d'_i}{d'_{\max}} \times 0.75
\end{equation}
giving a precision range of $[0.2, 0.95]$ with the empirical shape preserved. For the normal condition,
this yields a Beta$(5, 2)$ distribution (mean $\approx 0.64$, right-skewed), consistent with most bees
being above-chance at standard discrimination tasks.

\subsection*{Domain-specific observation functions}
\label{sec:model:domains}

Each behavioral domain maps precision to performance through a distinct functional form, reflecting
the different computational roles of the self-model.

\paragraph{Domain 1 --- Metacognition (opt-out accuracy).}
The opt-out task presents hard and easy stimuli. Correct metacognitive behavior requires opting out on
hard trials (where internal predictions are unreliable) and staying in on easy trials. The probability
of opting out on a trial with difficulty $d$ is:
\begin{equation}
  \poptout(d \mid \uncertainty) = \sigmoid\!\bigl(\uncertainty \cdot \slope - \thresh{d}\bigr)
\end{equation}
where $\sigmoid(\cdot)$ is the logistic function, $\slope = 4.0$ is the sensitivity, and $\thresh{d}$
is a difficulty-specific threshold ($\thresh{\text{hard}} = 1.0$, $\thresh{\text{easy}} = 2.5$).
Metacognitive accuracy is:
\begin{equation}
  \metaacc = \frac{\poptout(\text{hard} \mid \uncertainty) + (1 - \poptout(\text{easy} \mid \uncertainty))}{2}
  \label{eq:meta}
\end{equation}

\paragraph{Domain 2 --- Tool use anticipation.}
Tool use requires a forward model: the agent must anticipate the goal state before completing the
initial action step. Anticipatory orientation probability scales with precision weighted by caste
identity:
\begin{equation}
  P_{\text{tool}} = \precision \cdot \bigl(0.5 + 0.5 \cdot c\bigr)
\end{equation}
Foragers ($c \approx 1$) have higher goal-directedness, consistent with their spatial foraging role.

\paragraph{Domain 3 --- Caste-appropriate learning.}
Learning accuracy on a task of type $\tau$ depends on the match between task demands and caste
identity. For spatial tasks ($\tau = \text{spatial}$), foragers excel ($c \approx 1$); for social
tasks ($\tau = \text{social}$), nurses excel ($c \approx 0$). Formally:
\begin{equation}
  P_{\text{learn}}(\tau) = 0.5 + 0.45 \cdot \precision \cdot \text{match}(c, \tau)
\end{equation}
where $\text{match}(c, \text{spatial}) = c$ and $\text{match}(c, \text{social}) = 1 - c$.

\paragraph{Domain 4 --- General cognitive ability (GCA).}
GCA reflects the shared precision substrate across multiple tasks. We model a composite GCA score
as linear in precision:
\begin{equation}
  \text{GCA}_i = 0.5 + 0.4 \cdot \precision_i + \varepsilon_i, \quad \varepsilon_i \sim \mathcal{N}(0, 0.04)
\end{equation}

\subsection*{Conditions modeled}
\label{sec:model:conditions}

We simulate four experimental conditions (Table~\ref{tab:results}):
\begin{itemize}
  \item \textbf{Normal}: Beta$(5,2)$ precision distribution, mixed caste ($c \sim \mathcal{N}(0.5, 0.2)$)
  \item \textbf{Disrupted}: Beta$(2,4)$ precision distribution (circadian disruption reduces precision)
  \item \textbf{Nurse spatial}: Normal precision, nurse caste ($c \sim \mathcal{N}(0.15, 0.1)$), tested on spatial task
  \item \textbf{Forager spatial}: Slightly elevated precision, forager caste ($c \sim \mathcal{N}(0.85, 0.1)$), tested on spatial task
\end{itemize}
